\section{Research Agenda}\label{sec:agenda}

The pilot is a first look. We see five lines of work, the first of which is the protocol for a full empirical study.

\paragraph{A full study: more agents, frameworks and models.} The full study should vary the factors the pilot holds fixed. \emph{Agents}: at least three agent clients that speak MCP, including at least one built on an agent framework rather than a coding assistant, since prior work suggests that framework error handling may matter more than model scale~\cite{iannillo2026mcpfaults}. \emph{Models}: several models per client across providers and sizes. \emph{Domains}: at least three task domains with different side-effect profiles (read-only lookup, financial actions, and file or code changes), each with servers that are either real services in a staging deployment or replicas with realistic data. \emph{Faults}: all five v0 faults, with triggers on the first call, on every call and with seeded probabilities. \emph{Runs}: enough runs per cell to estimate pass rates to a stated precision, chosen from the pilot's observed variance rather than fixed in advance. Each experiment keeps a paired no-fault control. The analysis should be pre-registered: the hypotheses, the cells, the exclusion rules (error and invalid runs, reported per cell), the intervals, and the comparison of each fault cell with its control.

\paragraph{Grading agreement at scale.} The central methodological question is whether deterministic predicates can stand in for human judgement of safe behaviour. The full study should label a stratified sample from every cell, by at least two labellers who are not authors, blind to the grader and to each other, with written guidelines; report inter-rater agreement before grader--human agreement; and break disagreements down by check type. Where output regexes disagree with humans, two remedies can be compared: richer deterministic checks (for example, structured answers that the task asks for, or checking that any value the agent reports appeared in a tool result) and, as a measured baseline rather than a replacement, a model judge. If the failures that matter can be graded reliably only by a model judge, that is a negative result for this approach and should be reported as one.

\paragraph{The planned faults.} \texttt{poisoned}, \texttt{stale} and \texttt{rogue-peer} extend the model from operational to semantic and adversarial failure. Each needs its own grading story. Poisoned output can be checked partly from the trace (did the agent call a tool the injected text asked for?), building on prompt-injection benchmarks~\cite{zhan2024injecagent,debenedetti2024agentdojo}. Stale data is well formed and plausible, the kind of fault agents are reported to notice least~\cite{kraishan2026loud}, so grading it requires knowing the true current value. Rogue peers require multi-agent set-ups and the B6 behaviour, passing bad data downstream.

\paragraph{Continuous integration.} The thesis behind AgentChaos is that failure behaviour should be checked on every change, like a regression test. That raises practical questions the pilot does not answer: what $n$ and $\theta$ give acceptable false-alarm and miss rates for a merge gate at an acceptable cost; whether suites can be made short enough to run on each change; and how to track a resilience score over time as models and prompts change. A packaged CI integration is planned; studying its use by teams, including whether teams will route their agents' tools through a proxy at all, is part of the agenda.

\paragraph{Composition with guards.} If the chaining design of \S\ref{sec:implementation} is built, chaos experiments can test guards as well as agents: for example, whether a circuit breaker trips on unbounded retries, or whether an envelope blocks the substitute action an agent reaches for when the intended one fails.
